% Draft prepared 2026-10-11; elementary identities, no novelty claim.
% Intended to be copied into the coordinator's canonical claim record.
\begin{proposition}[An exact cubical intersection identity]
\label{prop:cubical-kinematics}
Let $n\ge1$, let $Q_n=\{0,1\}^n$, and let $S,T\subseteq Q_n$.
Write $\mathcal K(S)$ for the cubical subcomplex of $[0,1]^n$
consisting of every coordinate face all of whose vertices belong to $S$.
For $0\le j\le n$, let $f_j(S)$ be its number of $j$-dimensional faces;
in particular, $f_0(S)=|S|$. Set
\[
 N_j=2^{n-j}\binom nj,
 \qquad
 \chi(S)=\sum_{j=0}^n(-1)^j f_j(S).
\]
All these counts, including $\chi$, are zero when $S=\varnothing$.
Let $g$ be uniform in the finite group generated by coordinate permutations
and coordinate reflections $x_i\mapsto1-x_i$. Then
\begin{align}
 \mathbb E_g f_j(S\cap gT)&=\frac{f_j(S)f_j(T)}{N_j},
       &&0\le j\le n, \label{eq:cubical-face-kinematics}\\
 \mathbb E_g\chi(S\cap gT)&=
       \sum_{j=0}^n(-1)^j\frac{f_j(S)f_j(T)}{N_j}.
       \label{eq:cubical-euler-kinematics}
\end{align}
These are finite averaging identities; the Euler characteristic in the
second identity does not in general indicate nonemptiness.
\end{proposition}
\begin{proof}
The definition includes all subfaces of each admitted face, so it indeed
defines a cubical complex. A coordinate face belongs to $\mathcal K(S\cap
gT)$ precisely when it belongs to both $\mathcal K(S)$ and $g\mathcal K(T)$.
There are $N_j$ coordinate $j$-faces, and the specified group acts transitively
on them. For each fixed $j$-face $H$, the image $g^{-1}H$ is therefore
uniform among those $N_j$ faces. Thus
\[
 \Pr\{H\in g\mathcal K(T)\}=f_j(T)/N_j.
\]
Sum this equality over the $f_j(S)$ faces belonging to $\mathcal K(S)$.
This gives \eqref{eq:cubical-face-kinematics}; linearity of expectation
gives \eqref{eq:cubical-euler-kinematics}. The argument includes $j=0,n$
and the empty-set cases without modification.
\end{proof}

\begin{example}[SAT with zero Euler characteristic]
\label{ex:sat-cycle}
For $n=k=3$, the proper $3$-CNF
\[
 (x_1\lor x_2\lor x_3)\land
 (\neg x_1\lor\neg x_2\lor\neg x_3)
\]
has solution set $S=Q_3\setminus\{000,111\}$. Its induced cubical
complex has six vertices and six edges. Every square face contains one
of the two excluded vertices, so no square or three-dimensional face
is admitted. The complex is a six-cycle and
$\chi(S)=6-6=0$, although $S\ne\varnothing$.
Thus replacing the nonemptiness event by Euler characteristic in
\eqref{eq:cubical-euler-kinematics} is invalid even for a proper $3$-CNF
with two clauses.
\end{example}

\begin{example}[Face counts do not determine intersection probability]
\label{ex:face-profile-insufficient}
In $Q_3$, let $S=\{000,110\}$ and $U=\{000,111\}$.
Both sets have face vector $(2,0,0,0)$, but
\[
 \Pr_g\{S\cap gS\ne\varnothing\}=\frac5{12},
 \qquad
 \Pr_g\{S\cap gU\ne\varnothing\}=\frac12.
\]
Indeed, $gS$ is uniform among the twelve unordered distance-two pairs.
Exactly five such pairs meet $S$: each of its vertices belongs to three,
and their own pair was counted twice. In contrast, $gU$ is uniform among
the four antipodal pairs, of which exactly two meet $S$. Thus no function
of these two face vectors alone can equal the intersection probability
for all Boolean subsets. Each of $S$ and $U$ is a proper $3$-CNF solution
set, obtained by excluding its six missing vertices with six clauses.
\end{example}

\begin{proposition}[Expected face counts for random proper $k$-SAT]
\label{prop:expected-cubical-faces}
Let $n\ge k\ge2$ and $m\ge0$. Let $S_m$ be the solution set of
$F_k(n,m)$ in the project's independent-clause model. Define
\[
 q_j=\frac{1}{\binom nk}
       \sum_{r=0}^{k}\binom{n-j}{r}\binom{j}{k-r}2^{-r},
       \qquad 0\le j\le n,
\]
where out-of-range binomial coefficients are zero. Then, for $m\ge1$,
\begin{align}
 \mathbb E f_j(S_m)&=N_j(1-q_j)^m, \label{eq:random-cubical-faces}\\
 \mathbb E\chi(S_m)&=\sum_{j=0}^n(-1)^jN_j(1-q_j)^m.
\end{align}
For $m=0$, the expectations are $N_j$ and $1$, respectively.
In particular, $q_0=2^{-k}$ and $q_n=1$.
\end{proposition}
\begin{proof}
Fix a $j$-face, which has $j$ free coordinates and $n-j$ fixed coordinates.
The number $R$ of selected clause variables among the fixed coordinates
has probability
\[
 \Pr\{R=r\}=\binom{n-j}{r}\binom{j}{k-r}/\binom nk.
\]
The clause is falsified at some vertex of the face precisely when its
$R$ literals on fixed coordinates are all false. This has conditional
probability $2^{-R}$: the free coordinates can then be chosen to falsify
the remaining literals. Consequently the probability that every vertex
of the face satisfies one clause is $1-q_j$. Independence of the $m$
clauses gives $(1-q_j)^m$. Summing over all $N_j$ faces proves the first
identity, and alternating the finite sum proves the second. For $m=0$
all cube faces are present and the cube has Euler characteristic one.
\end{proof}

\begin{remark}[Positivity does not supply intrinsic-volume concentration]
The $f_j$ are nonnegative and invariant under the specified finite group.
They need not be log-concave. For example, in $Q_6$ let $S$ be the union
of the square $\{(a,b,0,0,0,0):a,b\in\{0,1\}\}$ and the $22$ vertices
of even total parity with at least two ones in their last four coordinates.
Those $22$ vertices are mutually nonadjacent and have distance at least
two from the square. Hence
\[
 (f_0(S),f_1(S),f_2(S),f_3(S),\ldots)=(26,4,1,0,\ldots),
 \qquad f_1(S)^2=16<26=f_0(S)f_2(S).
\]
The normalized sequence $f_j(S)/N_j$ also fails log-concavity at $j=1$.
This is an admissible proper $6$-CNF solution set: for each vertex outside
$S$, include the unique $6$-clause falsified at that vertex. This example
rules out a universal log-concavity assertion for these particular face
counts, not a more restricted typical-instance result or another invariant.
\end{remark}

\begin{proposition}[A coordinate-restriction Crofton identity]
\label{prop:coordinate-crofton}
Let $n\ge1$, $S\subseteq Q_n$, and $0\le r\le n$. Choose $I$ uniformly
among the $r$-element subsets of $\{1,\ldots,n\}$ and, conditional on $I$,
choose $b$ uniformly from $\{0,1\}^I$. Write
\[
 H_{I,b}=\{x\in Q_n:x_I=b\},
 \qquad \pi_I(S)=\{x_I:x\in S\}.
\]
Thus $H_{I,b}$ is the vertex set of a uniform coordinate face of
codimension $r$. Then
\begin{equation}
 \Pr\{S\cap H_{I,b}\ne\varnothing\}
 =\frac{1}{2^r\binom nr}\sum_{|I|=r}|\pi_I(S)|.
 \label{eq:coordinate-crofton}
\end{equation}
For $r=0$, the right-hand side is $\mathbf 1_{\{S\ne\varnothing\}}$;
for $r=n$, it is $|S|/2^n$. The identity is zero on both sides if
$S=\varnothing$.
\end{proposition}
\begin{proof}
For fixed $I$, the intersection is nonempty precisely for
$b\in\pi_I(S)$, so its conditional probability is $|\pi_I(S)|/2^r$.
Averaging over $I$ proves the formula. When $I=\varnothing$, the
projection of a nonempty set consists of the single empty tuple,
whereas the projection of the empty set is empty. This gives the
$r=0$ case; $r=n$ is the identity projection.
\end{proof}

\begin{proposition}[Intersection moments from ordered distances]
\label{prop:distance-intersection-moments}
Let $n\ge1$, let $S,T\subseteq Q_n$, and let $g$ be uniform in the
signed coordinate permutation group of
Proposition~\ref{prop:cubical-kinematics}. For $0\le d\le n$, define
\[
 A_d(S)=|\{(x,y)\in S\times S:\operatorname{dist}_H(x,y)=d\}|,
 \qquad Z=|S\cap gT|,
\]
where $\operatorname{dist}_H$ is Hamming distance and the pairs are ordered.
Then
\begin{align}
 \mathbb E Z&=\frac{|S||T|}{2^n},\label{eq:distance-first-moment}\\
 \mathbb E Z^2&=\sum_{d=0}^n
       \frac{A_d(S)A_d(T)}{2^n\binom nd}.
       \label{eq:distance-second-moment}
\end{align}
If both sets are nonempty, then $\mathbb E Z^2>0$ and
\begin{equation}
 \frac{(\mathbb E Z)^2}{\mathbb E Z^2}
 \le \Pr\{S\cap gT\ne\varnothing\}
 \le \min\{1,\mathbb E Z\}.
 \label{eq:distance-nonemptiness-bounds}
\end{equation}
If either set is empty, both moments and the intersection probability
are zero, and the quotient in \eqref{eq:distance-nonemptiness-bounds}
is not used.
\end{proposition}
\begin{proof}
Transitivity on vertices gives
$\Pr\{x\in gT\}=|T|/2^n$ for fixed $x$; summing over $x\in S$
proves \eqref{eq:distance-first-moment}.
There are $2^n\binom nd$ ordered pairs at distance $d$ in $Q_n$.
The group acts transitively on these pairs: translate the first member
to $0$, permute the $d$ differing coordinates to the target difference
set, and then translate $0$ to the target first member. These translations
are implemented by coordinate reflections. Consequently, for fixed $x,y$ at
distance $d$,
\[
 \Pr\{x,y\in gT\}
       =\frac{A_d(T)}{2^n\binom nd}.
\]
Expanding $Z^2$ as a sum of indicators over ordered pairs in $S^2$
and grouping by distance proves \eqref{eq:distance-second-moment}.
The argument includes $d=0$, where $A_0(S)=|S|$.
If both sets are nonempty, the first moment is positive and hence so
is the second. Finally, Cauchy--Schwarz gives
\[
 \mathbb E Z
 =\mathbb E[Z\mathbf 1_{\{Z>0\}}]
 \le (\mathbb E Z^2)^{1/2}\Pr\{Z>0\}^{1/2},
\]
which proves the lower bound. Since $Z$ is a nonnegative integer,
$\mathbf 1_{\{Z>0\}}\le Z$, proving the upper bound. Empty-set cases
follow directly from the definitions.
\end{proof}
